\scanpage{94}
\deleted{(à p. div. -- sous} Cat. dér. des Alg. de De Rham spéciales

Structure De Rham (à p.d. -- sous-entendu)

Anneau bigradué [à degrés -- total et ext -- $\geq 0$] alterné pour degré ext
\[
 A^{ij}\ne 0\quad\deleted{\unsure{\Leftarrow|et}}\Rightarrow
 \deleted{\unsure{j}}0\leq j\leq i
\]
degré complémentaire $=\unsure{j|i}-j\geq 0$

différentiel [différentielle de bidegré $0,1$ ; c'est une diff. pour le degré ext \deleted{\unsure{$j$|:}}
\[
d(xy)=(dx)y+(-1)^{\deg.\mathrm{ext}\,x}x\,dy
\]
]

à p. div. sur l'idéal $\displaystyle\sum_{i+j>0} A^{ij}=A^+$ [p. div. \deleted{\illegible{1}} compatible aux 2 grad \deleted{\unsure{$d$}} $(\gamma^d:A^{ij}\longrightarrow A^{di,dj})$ et la différentielle $(d\,x^{(i)}=x^{(i-1)}dx)$]

\uline{élément marqué} $T\in Z^{10}(A^{**})\simeq H^{10}(A^{**})$

[Cohomologie $H^{*i}(A^{**})\simeq\unsure{\tau_i|\prod_i|[}h^i\otimes_k k\{T\}$, \deleted{\unsure{où|$0$}}] $k=A^{00}$

Structure de De Rham \uline{spéciale} :
\[
\left\{\begin{aligned}
k\{T\}&\xrightarrow{\sim}H^{*0}(A^{**})\\
H^{*i}(A^{**})&=0\\
\text{les } A^{ij}&\text{ sont }k\text{-plats}
\end{aligned}\right.
\]
[\uline{NB} c'est utile pour la stabilité par $\otimes_k k'$]

(Structure de \unsure{D|De} Rham flasque : les $A^{ij}$ pour $j<i$ \inserted{sont flasques.}
% slanted
\margin{Univ. flasque : \deleted{\unsure{stabilité} \unsure{p}} \deleted{$\otimes_{\unsure{k}}$ \illegible{1}}}
% slanted

On parle de $k$-structure de De Rham si $k=A^{00}$ est donné.

\deleted{\illegible{1}} $\mathscr X$, topos

\deleted{\unsure{$\mathscr X$}} $\mathscr{DR}\mathrm{sp}(\mathscr X,k)_0$ structures de De Rham spéciales sur $\mathscr X$ -- avec comme flèches celles qui respectent \deleted{ttes} les structures \inserted{(y compris les structures)}

$\mathscr{DR}\mathrm{spfl}(\mathscr X,\deleted{\unsure{k}}k)_0$ \textquotedblleft\ spéciales et flasques

$\mathscr{DR}\mathrm{spufl}(\mathscr X,\deleted{\unsure{k}}k)_0$ \textquotedblleft\ \rule{1.5cm}{0.4pt} univ. flasques

$\mathscr{DR}\mathrm{sp}(\mathscr X,k)$, $\mathscr{DR}\mathrm{spfl}(\mathscr X,k)$, $\mathscr{DR}\mathrm{spufl}(\mathscr X,k)$

Catégories \inserted{dérivées} déduites en rendant inversibles ttes les flèches -- NB ce sont des quasi-isos --

% archive sheet number: 93
